% -------------------------------------------------------------------
% 2. Типографика~"--- определения и источники
% -------------------------------------------------------------------
\lecture[2026]{2}{Типографика}

\begin{attentionbox}{Что знать}
	Гарнитура, шрифт, начертание, кегль,
	классификация шрифтов (антиква, гротеск, акциденция), длина
	строки, интерлиньяж, трекинг, кернинг, «висячая»
	пунктуация, спецсимволы и~типографская раскладка, редактура
\end{attentionbox}

\begin{definition}[что такое «скучно»]
	<<Скучно>>~"--- лучший комментарий, который может получить начинающий
	дизайнер. «Скучно» означает, что в~вашем макете нет~явных ошибок,
	которые бросаются смотрящему в глаза и~портят впечатление
	о~макете. 
\end{definition}


\begin{codebox}[%
		Базовый HTML~"--- чтобы разбираться, как устроены веб-страницы
		с~помощью инспектора элементов%
	]
	\inputminted{html}{parts/code/html_1.html}
\end{codebox}


%\begin{tip}
%✦ Совет, полезная информация.
%\end{tip}

%\begin{solution}
%✦ Решение.
%\end{solution}

% vim:ts=2:sw=2:tw=68
